\chapter{CONCLUSÕES E TRABALHOS FUTUROS}

A partir do estudo acadêmico realizado e do trabalho dedicado a este sistema enquanto atividade laboral no LTI, conclui-se que o projeto de criação de um núcleo de informações centralizado no ambiente Linux foi extremamente produtivo e provou-se altamente útil para seus usuários.

Embora o uso inicial do protocolo gRPC para a comunicação da API voltada aos clientes finais não tenha trazido os benefícios de \textit{performance} esperados — face ao pequeno volume de dados (kilobytes) e ao "\textit{overhead}" de comunicação, sendo superado pela praticidade do protocolo HTTPS —, o aprendizado tecnológico obtido foi extremamente válido. Dessa forma, como trabalhos futuros, recomenda-se que esse conhecimento em gRPC seja aplicado em cenários arquiteturais adequados, tais como a comunicação interna de altíssimo rendimento e baixa latência entre dois ou mais microsserviços do próprio ecossistema no servidor (sem exposição à Internet).

Em relação ao impacto estrutural e organizacional, as atividades formais de engenharia de software introduzidas por este projeto evidenciaram à liderança que novas estruturas a nível de gerenciamento e governança não são apenas importantes, mas estritamente necessárias para a maturidade do LTI Finanças.

Por fim, os bons resultados de adoção demonstram que a estratégia adotada sanou uma dor real dos pesquisadores. Como próximos passos para trabalhos futuros, sugere-se a expansão dos conjuntos de dados integrados pela plataforma (novos ativos financeiros e resoluções temporais) e o refinamento contínuo das políticas de controle de acesso e processos de \textit{onboarding}, assegurando a escalabilidade e a manutenção do ecossistema a longo prazo.
